\section{Introduction}
A cluster operator sees a population of pending and running tasks. A workflow
analyst sees a population of unresolved applications. In both cases, a useful
question concerns their collective future: how many of the entities active now
will remain in their current state, and how uncertain is that count? Related
entities can evolve together because they belong to the same job, use the same
machine, or share a recent handler. A collection of accurate individual
probabilities does not specify this joint uncertainty.

Temporal graph learning provides useful representations of observed event
histories. However, an individual prediction benchmark and an operational
scenario benchmark answer different questions. Improving an entity's event
probability need not improve the distribution of a group count. Conversely,
adding shared random factors may improve a count distribution while leaving
every individual probability unchanged. Even an improvement in a proper
distributional score need not improve a particular asymmetric decision loss.
An application study should measure these distinctions directly.

We construct a common forecasting task on two public operational logs:
Google's 2011 cluster trace and BPI Challenge 2017. At a timestamp, an adapter
reconstructs the currently observed entities and typed relation memberships.
The target is their joint first-exit or resolution trajectory at four future
horizons. The entire current population is evaluated, including entities that
have not exited by the longest horizon. Future arrivals, eventual endpoints,
and future assignments are not supplied to the forecaster. This yields a
shared task across applications without assuming that a machine and a workflow
handler represent the same physical constraint.

Our study separates individual forecasts from the coupling of their outcomes.
We first compare tree predictors and temporal Mixer, GRU, and MLP ensembles.
We then hold the selected individual CDFs fixed while comparing independent,
global, typed relation, and context-dependent joint sampling. Classical
Gaussian and common-shock factors make this comparison computationally
feasible for changing populations. A risk-stratified incidence permutation
tests the role of observed relations. A historical aggregate-error bootstrap
provides a deliberately simple competitor that can correct group-level bias
but does not generate a consistent entity-level joint law.

The contribution is an application benchmark and controlled empirical study,
with reproducible adapters and a scalable reference implementation. It is not
a new copula family, a discovery that operational entities are dependent, or a
complete event-log simulator. The benchmark makes three questions separately
testable: whether temporal graph information improves individual forecasts;
whether observed relations improve collective forecasts when individual
probabilities are identical; and whether any distributional improvement
survives an operational loss and a future-period evaluation. Preserving negative
answers is part of the benchmark's purpose.
